\documentclass[a4paper]{article}
\usepackage[utf8]{inputenc} 	% codificacao de caracteres
\usepackage[T1]{fontenc}    	% codificacao de fontes
\usepackage[brazil]{babel}  	% idioma
\usepackage[table]{xcolor}    	% linhas coloridas alternadamente
\usepackage{graphicx}
\usepackage{multirow}
\usepackage{soulutf8}			% cor de fundo do texto que aceita utf8
\usepackage{makecell}			% mais de uma linha na célula da tabela

% espaçamento paragrafos
\setlength{\parindent}{0em}
%\setlength{\parskip}{2em}
%\renewcommand{\baselinestretch}{1}

%opening
\title{Universidade Federal do Ceará - Crateús}
\author{Ítalo Ribeiro}
\date{}

\begin{document}

\fontfamily{cmss}
\fontsize{12}{18}
\selectfont

\begin{center}
	\includegraphics[scale=0.18]{brasao2_vertical}
\end{center}

\begin{center}
	\vspace{2ex}
	Campus de Crateús\\	
	Estágio Supervisionado\\
	\vspace{4ex}
	\textbf{{
			\large PLANO DE TRABALHO DE ESTÁGIO SUPERVISIONADO
	}}\\
	\vspace{2ex}
	{\large 
		Estágio <Especificar Instituição>\\
		<Título do Projeto que está sendo desenvolvido>
	}
\end{center}

\section*{Identificação}

\textbf{Estagiário:}\\

Nome:  \\
Matrícula:  \\
Telefones:  \\
E-mail:  \\
Cargo:  \\

\textbf{Instituição do Estágio:}\\

Nome:  \\
Endereço:  \\
Ramo de Atividade:  \\
Supervisor:  \\
Telefone:  \\
E-mail:  \\

\textbf{Professor Orientador:}

\newpage

\section*{Introdução}

\sethlcolor{lightgray}

\textit{\hl{[Nesta seção deve ser falado sobre a instituição que você desenvolvendo o estágio e sobre as atividades que serão desenvolvidas em linhas gerais (Se for um sistema, falar das características gerais e se as atividades estão ligadas ao desenvolvimento, análise, etc.).]}}

\section*{Qualificação do aluno}

\textit{\hl{[Nesta seção você deve inserir as aptidões que tem, conhecimentos,
experiências profissionais, publicações (caso houver).]}}


\section*{Objetivos Gerais do Estágio}

\begin{itemize}
	\item \textit{\hl{Objetivo 1 do estágio;}}
	\item \textit{\hl{Objetivo 2 do estágio;}}
	\item \textit{\hl{Objetivo 3 do estágio;}}
	\item \textit{\hl{Etc.}}
\end{itemize}

\section*{Justificativa do Tema}

\textit{\hl{[Justificar a importância do local de trabalho e das atividades a serem realizadas para sua carreira profissional.]}}


\section*{Cronograma de Atividades}

\textit{\hl{[Elencar as possíveis atividades e a duração de acordo com o mês.]}}\\

\begin{tabular}{|c|c|c|c|c|c|c|c|}
	\hline
	\textbf{Atividade/Mês} &  &  &  &  &  &  &  \\
	\hline
	\makecell[l]{Se a celula\\ da tabela\\ possuir mais\\ de uma linha}&  &  &  &  &  &  &  \\
	\hline
\end{tabular}

\end{document}
